% 书籍第 4 册（《大学物理》第二年）中照片与 AI 插图的图片来源。机器可读的许可记录保存在 images/book4/CREDITS.md 中；两个文件需保持同步。
\chapter*{图片来源}

本书中的绘图均为原创。照片依其各自的自由许可使用，谨向作者致谢：

\begin{itemize}
  \item \emph{水电站压力管道}（第 5 章）：Chinchu.c，Wikimedia Commons，CC0。
  \item \emph{云中的卡门涡街}（第 4 章）：Robert Cahalan，NASA/GSFC，公有领域。
  \item \emph{同轴电缆}（第 8 章）：Mx.~Granger，Wikimedia Commons，CC0。
  \item \emph{詹姆斯·克拉克·麦克斯韦}（第 11 章）：Smithsonian Institution，公有领域。
  \item \emph{国际空间站拍摄的南极光}（第 14 章）：Joe Acaba，NASA，公有领域。
  \item \emph{光学平晶干涉条纹}（第 19 章）：\emph{Machinery} 杂志，约 1915 年，公有领域。
  \item \emph{LIGO 汉福德站航拍图}（第 20 章）：LIGO Laboratory，公有领域。
  \item \emph{594~nm 氦氖激光器}（第 23 章）：Telementor，Wikimedia Commons，CC~BY~4.0。
  \item \emph{极紫外波段的太阳}（第 26 章）：NASA/SDO (AIA)，公有领域。
  \item \emph{阿波罗 17 号拍摄的地球}（第 26 章）：NASA/Apollo 17 机组，公有领域。
  \item \emph{路德维希·玻尔兹曼}（第 29 章）：摄影者不详，公有领域。
  \item \emph{电子衍射图样}（第 30 章）：National Institute of Standards and Technology，公有领域。
  \item \emph{STM 拍摄的金（100）表面}（第 31 章）：Erwin Rossen，Wikimedia Commons，公有领域。
  \item \emph{查尔斯·汤斯与第一台微波激射器}（第 31 章）：Dan Rubin，公有领域。
\end{itemize}

每张照片的完整来源链接与许可 URL 列于本书代码库中的 \texttt{images/book4/CREDITS.md}。

\bigskip
除照片外，第 1--10 章和第 12--31 章中的三十幅图为 AI 生成插图，由 OpenAI 图像生成功能依据本书编辑撰写的提示词生成，并在收录前核对其物理准确性。每幅生成图像的完整提示词记录于代码库中的 \texttt{images/book4/ai/PROMPTS.md}。
\clearpage
